\ifdefined\gkver\else\def\gkver{student}\fi
\documentclass[\gkver]{gaokaozhenti}
\begin{document}

\chapter{1989年全国}

\begin{enumerate}
\item
%% number: 1
%% paperName: 1989年全国高考第 1 题 2 分
%% typeId: 1
%% type: 单选题
%% chapter: 曲线运动
%% point: 万有引力定律
%% method: 无
%% score: 2
%% degree: 850
%% duplicateId: 0
%% body:
设地球表面的重力加速度为 $g_{0}$，物体在距地心 $4R$（$R$ 是地球半径）处，由于地球的作用而产生的加速度为 $g$，则 $\frac{g}{g_{0}}$ 为\xzanswer{D}
\fourchoices[answer=D]
{1}
{1/9}
{1/4}
{1/16}

%% answer: D
%% memo:
\memoanswer{
地球上的重力加速度由万有引力产生，有 $G\frac{Mm}{R^{2}}=mg_{0}$，$G\frac{Mm}{(4R)^{2}}=mg$，联立解得 $\frac{g}{g_{0}}=\frac{1}{16}$，故 D 正确。
}

\item
%% number: 2
%% paperName: 1989年全国高考第 2 题 2 分
%% typeId: 1
%% type: 单选题
%% chapter: 机械振动
%% point: 单摆
%% method: 无
%% score: 2
%% degree: 850
%% duplicateId: 0
%% body:
若单摆的摆长不变，摆球的质量增加为原来的 4 倍，摆球经过平衡位置时的速度减小为原来的 1/2，则单摆振动的\xzanswer{B}
\fourchoices[answer=B]
{频率不变，振幅不变}
{频率不变，振幅改变}
{频率改变，振幅改变}
{频率改变，振幅不变}

%% answer: B
%% memo:
\memoanswer{
由单摆周期公式知 $T=2\pi\sqrt{\frac{l}{g}}$，$l$ 不变，$T$ 不变，又 $f=\frac{1}{T}$，即 $f$ 不变；由机械能守恒得 $\frac{1}{2}mv^{2}=\frac{1}{2}kA^{2}$，$\frac{1}{2}m\left(\frac{v}{2}\right)^{2}=\frac{1}{2}kA'^{2}$，解得 $A'=\frac{A}{2}$，振幅改变，故 B 正确。
}

\item
%% number: 3
%% paperName: 1989年全国高考第 3 题 2 分
%% typeId: 1
%% type: 单选题
%% chapter: 磁场
%% point: 带电粒子在磁场中的运动
%% method: 无
%% score: 2
%% degree: 650
%% duplicateId: 0
%% body:
有三束粒子，分别是质子（p）、氚核（$\ce{^{3}_{1}H}$）和 $\alpha$ 粒子束，如果它们以相同的速度沿垂直于磁场方向射入匀强磁场（磁场方向垂直纸面向里）。在下列四图中，哪个图正确地表示出这三束粒子的运动轨迹？\xzanswer{C}
\fourchoices[answer=C, ispicture=true, h=2.2cm]
{figs/03a.png}
{figs/03b.png}
{figs/03c.png}
{figs/03d.png}

%% answer: C
%% memo:
\memoanswer{
质子 p 为 $\ce{^{1}_{1}H}$，氚核为 $\ce{^{3}_{1}H}$，$\alpha$ 粒子为 $\ce{^{4}_{2}He}$，又 $R=\frac{mv}{qB}$，所以 $R_{p}:R_{\text{氚}}:R_{\alpha}=\frac{1}{1}:\frac{3}{1}:\frac{4}{2}=1:3:2$，故 C 正确。
}

\item
%% number: 4
%% paperName: 1989年全国高考第 4 题 2 分
%% typeId: 1
%% type: 单选题
%% chapter: 运动和力
%% point: 牛顿定律的应用
%% method: 无
%% score: 2
%% degree: 750
%% duplicateId: 0
%% body:
两辆汽车在同一平直路面上行驶，它们的质量之比 $m_{1}:m_{2}=1:2$，速度之比 $v_{1}:v_{2}=2:1$。当两车急刹车后，甲车滑行的最大距离为 $s_{1}$，乙车滑行的最大距离为 $s_{2}$。设两车与路面间的滑动摩擦系数相等，不计空气阻力，则\xzanswer{D}
\fourchoices[answer=D]
{$s_{1}:s_{2}=1:2$}
{$s_{1}:s_{2}=1:1$}
{$s_{1}:s_{2}=2:1$}
{$s_{1}:s_{2}=4:1$}

%% answer: D
%% memo:
\memoanswer{
对两辆汽车运动过程，由动能定理得 $\mu mgs=\frac{1}{2}mv^{2}$，解得 $s=\frac{v^{2}}{2\mu g}$，得 $\frac{s_{1}}{s_{2}}=\left(\frac{v_{1}}{v_{2}}\right)^{2}=\frac{4}{1}$，故 D 正确。
}

\item
%% number: 5
%% paperName: 1989年全国高考第 5 题 2 分
%% typeId: 1
%% type: 单选题
%% chapter: 运动和力
%% point: 牛顿第二定律
%% method: 无
%% score: 2
%% degree: 750
%% duplicateId: 0
%% body:
一轻弹簧上端固定，下端挂一重物，平衡时弹簧伸长了 $4\Ucm$。再将重物向下拉 $1\Ucm$，然后放手，则在刚释放的瞬间重物的加速度是（$g$ 取 $10\Umsq$）\xzanswer{A}
\fourchoices[answer=A]
{$2.5\Umsq$}
{$7.5\Umsq$}
{$10\Umsq$}
{$12.5\Umsq$}

%% answer: A
%% memo:
\memoanswer{
重物平衡时，有 $mg=k\times0.04$；刚释放时，有 $k\times0.05-mg=ma$，解得 $a=\frac{1}{4}g=2.5\Umsq$，故 A 正确。
}

\item
%% number: 6
%% paperName: 1989年全国高考第 6 题 2 分
%% typeId: 1
%% type: 单选题
%% chapter: 功和能
%% point: 动能定理
%% method: 无
%% score: 2
%% degree: 650
%% duplicateId: 0
%% body:
一质量为 $m$ 的小球，用长为 $l$ 的轻绳悬挂于 O 点。小球在水平拉力 $F$ 作用下，从平衡位置 P 点很缓慢地移动到 Q 点（如图所示），则力 $F$ 所做的功为\xzanswer{B}
\onepicture[w=3cm]{\includesvg{figs/06}}
\fourchoices[answer=B]
{$mgl\cos\theta$}
{$mgl(1-\cos\theta)$}
{$Fl\sin\theta$}
{$Fl\theta$}

%% answer: B
%% memo:
\memoanswer{
对小球的运动过程，由动能定理得 $W_{F}-mgl(1-\cos\theta)=0$，解得 $W_{F}=mgl(1-\cos\theta)$，故 B 正确。
}

\item
%% number: 7
%% paperName: 1989年全国高考第 7 题 2 分
%% typeId: 1
%% type: 单选题
%% chapter: 电路
%% point: 电容
%% method: 无
%% score: 2
%% degree: 650
%% duplicateId: 0
%% body:
在图示的电路中，已知电容 $C=2\UuF$，电源电动势 $\varepsilon=12\UV$，内电阻不计，$R_{1}:R_{2}:R_{3}:R_{4}=1:2:6:3$。则电容器极板 a 所带的电量为\xzanswer{D}
\onepicture[w=3.5cm]{figs/07.png}
\fourchoices[answer=D]
{$-8\times10^{-6}\UC$}
{$4\times10^{-6}\UC$}
{$-4\times10^{-6}\UC$}
{$8\times10^{-6}\UC$}

%% answer: D
%% memo:
\memoanswer{
由串联电路的分压特性知，$R_{1}$、$R_{2}$ 分压，$U_{a}=E\times\frac{R_{2}}{R_{1}+R_{2}}=8\UV$；$R_{3}$、$R_{4}$ 分压，$U_{b}=E\times\frac{R_{4}}{R_{3}+R_{4}}=4\UV$，故 $U_{ab}=U_{a}-U_{b}=4\UV$，又 $Q=CU_{ab}=8\times10^{-6}\UC$，故 D 正确。
}

\item
%% number: 8
%% paperName: 1989年全国高考第 8 题 2 分
%% typeId: 1
%% type: 单选题
%% chapter: 电路
%% point: 动态分析
%% method: 无
%% score: 2
%% degree: 750
%% duplicateId: 0
%% body:
在如图所示的电路中，当滑线变阻器的滑动触点向 b 端移动时，\xzanswer{A}
\onepicture[w=4cm]{figs/08.png}
\fourchoices[answer=A]
{伏特表 V 的读数增大，安培表 A 的读数减小}
{伏特表 V 和安培表 A 的读数都增大}
{伏特表 V 和安培表 A 的读数都减小}
{伏特表 V 的读数减小，安培表 A 的读数增大}

%% answer: A
%% memo:
\memoanswer{
滑动触点向 b 移动，$R_{3}$ 变大，干路电流 $I$ 变小，由 $U=E-Ir$ 知，电压表 V 的读数 $U$ 增大，电流表 A 的读数 $I_{3}=I-I_{2}$，因 $I_{2}=\frac{U}{R_{2}}=\frac{E-I(r+R_{1})}{R_{2}}$ 变大，所以 $I_{3}$ 变小，故 A 正确。
}

\item
%% number: 9
%% paperName: 1989年全国高考第 9 题 2 分
%% typeId: 1
%% type: 单选题
%% chapter: 磁场
%% point: 左手定则
%% method: 无
%% score: 2
%% degree: 650
%% duplicateId: 0
%% body:
一矩形通电线框 abcd，可绕其中心轴 $OO'$ 转动，它处在与 $OO'$ 垂直的匀强磁场中（如图）。在磁场作用下线框开始转动，最后静止在平衡位置。则平衡后\xzanswer{B}
\onepicture[w=3cm]{figs/09.png}
\fourchoices[answer=B]
{线框四边都不受磁场的作用力}
{线框四边受到指向线框外部的磁场作用力，但合力为零}
{线框四边受到指向线框内部的磁场作用力，但合力为零}
{线框的一对边受到指向线框外部的磁场作用力，另一对边受到指向线框内部的磁场作用力，但合力为零}

%% answer: B
%% memo:
\memoanswer{
由左手定则知，线框逆时针转动（俯视），转到中性面（线框垂直于 B）静止，静止时，由左手定则知，受力分析如图所示。

\onepicture[w=3cm]{figs/09_an.png}

故 B 正确。
}

\item
%% number: 10
%% paperName: 1989年全国高考第 10 题 2 分
%% typeId: 1
%% type: 单选题
%% chapter: 曲线运动
%% point: 平抛运动
%% method: 无
%% score: 2
%% degree: 650
%% duplicateId: 0
%% body:
一架飞机水平地匀速飞行。从飞机上每隔 1 秒钟释放一个铁球，先后共释放 4 个。若不计空气阻力，则四个球\xzanswer{C}
\fourchoices[answer=C]
{在空中任何时刻总是排成抛物线；它们的落地点是等间距的}
{在空中任何时刻总是排成抛物线；它们的落地点是不等间距的}
{在空中任何时刻总在飞机正下方排成竖直的直线；它们的落地点是等间距的}
{在空中任何时刻总在飞机正下方排成竖直的直线；它们的落地点是不等间距的}

%% answer: C
%% memo:
\memoanswer{
平抛运动水平方向做匀速运动，且铁球的水平速度与飞机相同，每个小球释放后，小球均做抛物线运动；以飞机为参考系，铁球做自由落体运动，落地的时间差与释放的时间差相同，所以落地点等距，故 C 正确。
}

\item
%% number: 11
%% paperName: 1989年全国高考第 11 题 2 分
%% typeId: 1
%% type: 单选题
%% chapter: 气体性质
%% point: 气体图象
%% method: 无
%% score: 2
%% degree: 650
%% duplicateId: 0
%% body:
$p-T$ 图上的图线 abc 表示一定质量的理想气体的状态变化过程，此过程在 $p-V$ 图上的图线应为\xzanswer{C}
\onepicture[w=3.5cm]{\includesvg{figs/11}}
\fourchoices[answer=C, ispicture=true, h=2.2cm]
{figs/11a.png}
{figs/11b.png}
{figs/11c.png}
{figs/11d.png}

%% answer: C
%% memo:
\memoanswer{
分析 $p-T$ 图像，由图象知 $a\to b$ 为等容过程，且 $p_{a}<p_{b}$；$b\to c$ 为等温过程，由等温过程得 $p_{c}V_{c}=p_{b}V_{b}$，又 $p_{b}>p_{c}$，故 $V_{b}<V_{c}$，所以在 $p-V$ 图上，$a\to b$ 图象垂直于 V 轴，$b\to c$ 图象为双曲线（反比），故 C 正确。
}

\item
%% number: 12
%% paperName: 1989年全国高考第 12 题 2 分
%% typeId: 1
%% type: 单选题
%% chapter: 电磁感应
%% point: 自感与互感，涡流
%% method: 无
%% score: 2
%% degree: 750
%% duplicateId: 0
%% body:
在下列四个日光灯的接线图中（S 为起辉器，L 为镇流器），正确的是\xzanswer{A}
\fourchoices[answer=A, ispicture=true, h=2.4cm]
{figs/12a.png}
{figs/12b.png}
{figs/12c.png}
{figs/12d.png}

%% answer: A
%% memo:
\memoanswer{
由日光灯的发光原理知，通过启辉器的开闭，使镇流器两端的电压达到日光灯的启动电压，电路稳定后，镇流器应串联在主电路中，起到稳压限流的作用，故 A 正确。
}

\item
%% number: 13
%% paperName: 1989年全国高考第 13 题 2 分
%% typeId: 1
%% type: 单选题
%% chapter: 运动和力
%% point: 动量守恒定律
%% method: 无
%% score: 2
%% degree: 650
%% duplicateId: 0
%% body:
在光滑水平面上有三个完全相同的小球排成一条直线。2、3 小球静止，并靠在一起，1 球以速度 $v_{0}$ 射向它们（如图）。设碰撞中不损失机械能，则碰后三个小球的速度可能值是\xzanswer{D}
\onepicture[w=5cm]{\includesvg{figs/13}}
\fourchoices[answer=D]
{$v_{1}=v_{2}=v_{3}=\frac{1}{\sqrt{3}}v_{0}$}
{$v_{1}=0$，$v_{2}=v_{3}=\frac{1}{\sqrt{2}}v_{0}$}
{$v_{1}=0$，$v_{2}=v_{3}=\frac{1}{2}v_{0}$}
{$v_{1}=v_{2}=0$，$v_{3}=v_{0}$}

%% answer: D
%% memo:
\memoanswer{
由题意知，碰撞过程中满足动量守恒和机械能守恒。由动量守恒得 AB 错误，CD 可能正确；C 项：初始时系统的动能为 $\frac{1}{2}mv_{0}^{2}$，碰后系统的动能为 $2\times\frac{1}{2}m\left(\frac{1}{2}v_{0}\right)^{2}=\frac{1}{4}mv_{0}^{2}$，知系统初末动能不守恒，故 C 错误；D 项：系统的初末动能守恒，故 D 正确。
}

\item
%% number: 14
%% paperName: 1989年全国高考第 14 题 2 分
%% typeId: 1
%% type: 单选题
%% chapter: 光学
%% point: 透镜
%% method: 无
%% score: 2
%% degree: 650
%% duplicateId: 0
%% body:
红色、绿色和黄色的三束平行光分别沿主轴射向同一个玻璃凸透镜，通过透镜后会聚到主轴上，会聚点到光心的距离分别是 $f_{\text{红}}$、$f_{\text{绿}}$、$f_{\text{黄}}$，则\xzanswer{C}
\fourchoices[answer=C]
{$f_{\text{红}}=f_{\text{绿}}=f_{\text{黄}}$}
{$f_{\text{红}}<f_{\text{黄}}<f_{\text{绿}}$}
{$f_{\text{绿}}<f_{\text{黄}}<f_{\text{红}}$}
{$f_{\text{红}}>f_{\text{绿}}>f_{\text{黄}}$}

%% answer: C
%% memo:
\memoanswer{
将三种光按照频率从小到大排序：红、黄、绿，光的频率越大，折射率越大，偏折越厉害，焦距越短，故 C 正确。
}

\item
%% number: 15
%% paperName: 1989年全国高考第 15 题 2 分
%% typeId: 1
%% type: 单选题
%% chapter: 力物体的平衡
%% point: 力矩平衡
%% method: 无
%% score: 2
%% degree: 650
%% duplicateId: 0
%% body:
在光滑水平地面上有一木板，一木棒可沿水平轴 O 转动，其下端 B 搁在木板上，而整个系统处于静止状态（如图）。现在用水平力 $F$ 向左推木板，但木板仍未动。由此可以得出结论：施力 $F$ 后，木板和木棒之间的正压力\xzanswer{C}
\onepicture[w=4.5cm]{\includesvg{figs/15}}
\fourchoices[answer=C]
{变大}
{不变}
{变小}
{条件不足，不能判断如何改变}

%% answer: C
%% memo:
\memoanswer{
设棒长为 $L$，由题意知，初始时，棒与木板间无摩擦力的作用，施力后棒与木板间有摩擦力的作用，由力矩平衡得，施力前 $G\frac{L}{2}=NL$，施力后 $G\frac{L}{2}=N'L+FL$，所以 $N'<N$，故 C 正确。
}

\item
%% number: 16
%% paperName: 1989年全国高考第 16 题 2 分
%% typeId: 2
%% type: 多选题
%% chapter: 原子和原子核
%% point: 玻尔模型
%% method: 无
%% score: 2
%% degree: 750
%% duplicateId: 0
%% body:
玻尔在他提出的原子模型中所做的假设有\xzanswer{ABC}
\fourchoices[answer=ABC]
{原子处于称为定态的能量状态时，虽然电子做加速运动，但并不向外辐射能量}
{原子的不同能量状态与电子沿不同的圆轨道绕核运动相对应，而电子的可能轨道的分布是不连续的}
{电子从一个轨道跃迁到另一轨道时，辐射（或吸收）一定频率的光子}
{电子跃迁时辐射的光子的频率等于电子绕核做圆周运动的频率}

%% answer: ABC
%% memo:
\memoanswer{
由玻尔假设的内容知 ABC 正确。
}

\item
%% number: 17
%% paperName: 1989年全国高考第 17 题 2 分
%% typeId: 1
%% type: 单选题
%% chapter: 电场
%% point: 等势面
%% method: 无
%% score: 2
%% degree: 750
%% duplicateId: 0
%% body:
一个点电荷，从静电场中的 a 点移到 b 点，其电势能的变化为零，则\xzanswer{D}
\fourchoices[answer=D]
{a、b 两点的场强一定相等}
{该点电荷一定沿等势面移动}
{作用于该点电荷的电场力与其移动方向总是垂直的}
{a、b 两点的电势一定相等}

%% answer: D
%% memo:
\memoanswer{
由功能关系得 $W=\Delta E=q(U_{a}-U_{b})=0$，所以 $U_{a}=U_{b}$，故 D 正确。
}

\item
%% number: 18
%% paperName: 1989年全国高考第 18 题 2 分
%% typeId: 2
%% type: 多选题
%% chapter: 电路
%% point: 动态分析
%% method: 无
%% score: 2
%% degree: 550
%% duplicateId: 0
%% body:
一平行板电容器 C，极板是水平放置的，它和三个可变电阻及电源联接成如图所示的电路。今有一质量为 $m$ 的带电油滴悬浮在两极板之间静止不动。要使油滴上升，可采用的办法是\xzanswer{CD}
\onepicture[w=4.5cm]{figs/18.png}
\fourchoices[answer=CD]
{增大 $R_{1}$}
{增大 $R_{2}$}
{增大 $R_{3}$}
{减小 $R_{2}$}

%% answer: CD
%% memo:
\memoanswer{
由串联电路的分压原理知，$R_{3}$ 的电压为 $U=\frac{ER_{3}}{R_{2}+R_{3}}$，要增大 $U$，须增大 $R_{3}$ 或减小 $R_{2}$，故 CD 正确。注意：$R_{1}$ 无电流通过，此时 $R_{1}$ 相当于导线。
}

\item
%% number: 19
%% paperName: 1989年全国高考第 19 题 2 分
%% typeId: 2
%% type: 多选题
%% chapter: 电路
%% point: 电容
%% method: 无
%% score: 2
%% degree: 550
%% duplicateId: 0
%% body:
一平行板电容器充电后与电源断开，负极板接地。在两极板间有一正电荷（电量很小）固定在 P 点，如下图所示。以 $E$ 表示两极板间的场强，$U$ 表示电容器的电压，$W$ 表示正电荷在 P 点的电势能。若保持负极板不动，将正极板移到图中虚线所示的位置，则\xzanswer{AC}
\onepicture[w=3.5cm]{figs/19.png}
\fourchoices[answer=AC]
{$U$ 变小，$E$ 不变}
{$E$ 变大，$W$ 变大}
{$U$ 变小，$W$ 不变}
{$U$ 不变，$W$ 不变}

%% answer: AC
%% memo:
\memoanswer{
由题意知电容两端的 $Q$ 保持不变，$d$ 减小，$C$ 增大，又 $U=\frac{Q}{C}$，故 $U$ 变小；$E=\frac{U}{d}=\frac{Q}{C\cdot d}$，又 $C=\frac{\varepsilon S}{4\pi kd}$，得 $E=\frac{Q}{d}\cdot\frac{4\pi kd}{\varepsilon S}=\frac{4\pi kQ}{\varepsilon S}$，故 $E$ 不变；下极板接地，设 P 点的电势为 $\varphi_{P}$，$E_{\text{p}}=q(\varphi_{P}-0)=qU_{P0}=qEd_{P}$，$d_{P}$ 不变，所以 $E_{\text{p}}$ 不变，故 AC 正确。
}

\item
%% number: 20
%% paperName: 1989年全国高考第 20 题 2 分
%% typeId: 2
%% type: 多选题
%% chapter: 气体性质
%% point: 热力学第一定律
%% method: 无
%% score: 2
%% degree: 450
%% duplicateId: 0
%% body:
对于一定质量的理想气体，在下列各种过程中，可能发生的过程是\xzanswer{ABCD}
\fourchoices[answer=ABCD]
{气体膨胀对外做功，温度升高}
{气体吸热，温度降低}
{气体放热，压强增大}
{气体放热，温度不变}

%% answer: ABCD
%% memo:
\memoanswer{
由热力学第一定律知 $\Delta U=W+Q$，当 $W$ 为负时，若 $Q$ 为正，且大于 $W$，则 $\Delta U$ 为正，故 A 正确；当 $Q$ 为正时，若 $W$ 为负，且绝对值大于 $Q$，则 $\Delta U$ 为负，故 B 正确；当 $Q$ 为负时，若 $W$ 为正，且大于 $Q$，则 $\Delta U$ 为正，故 C 正确；当 $Q$ 为负时，若 $W$ 为正，且 $W=|Q|$，则 $\Delta U=0$，故 D 正确。
}

\item
%% number: 21
%% paperName: 1989年全国高考第 21 题 3 分
%% typeId: 3
%% type: 填空题
%% chapter: 原子和原子核
%% point: 结合能
%% method: 无
%% score: 3
%% degree: 850
%% duplicateId: 0
%% body:
中子的质量为 $1.0087\Uu$，质子的质量为 $1.0073\Uu$，氘核的质量为 $2.0136\Uu$。中子和质子结合成氘核时释放的能量为\tkanswer[2.6cm]{$3.7\times10^{-13}$} 焦耳。（计算结果取两位有效数字。$1\Uu=1.7\times10^{-27}$ 千克）

%% answer: 3.7\times10^{-13}
\jdanswer{
$3.7\times10^{-13}$
}

%% memo:
\memoanswer{
该核反应方程为 $\ce{^{1}_{1}H + ^{1}_{0}n -> ^{2}_{1}H}$，质量亏损为 $\Delta m=(1.0087+1.0073-2.0136)\Uu=0.0024\Uu$，由质能方程得 $\Delta E=\Delta m\cdot c^{2}=3.7\times10^{-13}\UJ$。
}

\item
%% number: 22
%% paperName: 1989年全国高考第 22 题 3 分
%% typeId: 3
%% type: 填空题
%% chapter: 气体性质
%% point: 气体连接体
%% method: 无
%% score: 3
%% degree: 550
%% duplicateId: 0
%% body:
图中 A、B 是体积相同的气缸，B 内有一导热的、可在气缸内无摩擦滑动的、体积不计的活塞 C，D 为不导热的阀门。起初，阀门关闭，A 内装有压强 $p_{1}=2.0\times10^{5}\UPa$，温度 $T_{1}=300\UK$ 的氮气。B 内装有压强 $p_{2}=1.0\times10^{5}\UPa$，温度 $T_{2}=600\UK$ 的氧气。阀门打开后，活塞 C 向右移动，最后达到平衡。以 $V_{1}$ 和 $V_{2}$ 分别表示平衡后氮气和氧气的体积，则 $V_{1}:V_{2}=$\tkanswer[1.2cm]{4}。（假定氧气和氮气均为理想气体，并与外界无热交换，连接气缸的管道体积可忽略）
\onepicture[align=right, w=5cm]{figs/22.png}

%% answer: 4
\jdanswer{
4
}

%% memo:
\memoanswer{
系统平衡后两侧气体压强相等为 $p$，温度相等为 $T$，由理想气体状态方程得 $\frac{p_{1}V}{T_{1}}=\frac{pV_{1}}{T}$，$\frac{p_{2}V}{T_{2}}=\frac{pV_{2}}{T}$，联立化简得 $\frac{V_{1}}{V_{2}}=\frac{p_{1}T_{2}}{p_{2}T_{1}}=\frac{4}{1}$，由几何关系得 $V_{1}+V_{2}=2V$，解得 $V_{1}=\frac{8}{5}V$，$V_{2}=\frac{2}{5}V$。
}

\item
%% number: 23
%% paperName: 1989年全国高考第 23 题 3 分
%% typeId: 3
%% type: 填空题
%% chapter: 机械波
%% point: 波的多解性
%% method: 无
%% score: 3
%% degree: 550
%% duplicateId: 0
%% body:
下图中实线是一列简谐波在某一时刻的波形图线。虚线是 $0.2\Us$ 后它的波形图线。这列波可能的传播速度是\tkanswer[6cm]{$20\left(n+\frac{1}{4}\right)\Ums$ 或 $20\left(n+\frac{3}{4}\right)\Ums$，$n$ 取自然数}。
\onepicture[align=right, w=8cm]{\includesvg{figs/23}}

%% answer: 正向传播：$v=20\left(n+\frac{1}{4}\right)\Ums$，$n=0,1,2\cdots\cdots$；反向传播：$v=20\left(n+\frac{3}{4}\right)\Ums$，$n=0,1,2\cdots\cdots$
\jdanswer{
\begin{enumerate}
	\item
	正向传播：$v=20\left(n+\frac{1}{4}\right)\Ums$，$n=0,1,2\cdots\cdots$
	\item
	反向传播：$v=20\left(n+\frac{3}{4}\right)\Ums$，$n=0,1,2\cdots\cdots$
\end{enumerate}
}

%% memo:
\memoanswer{
由题图知，波长 $\lambda=4\Um$，波传播的时间为 $\Delta t=0.2\Us$。若波沿 $x$ 轴正向传播，由题图知波传播的最短距离为 $\frac{\lambda}{4}$，即传播的最短时间为 $\frac{T}{4}$，由波的周期性知 $\Delta t=\left(n+\frac{1}{4}\right)T$，$n$ 取自然数，由 $v=\frac{\lambda}{T}$ 得 $v=\frac{4\left(n+\frac{1}{4}\right)}{\Delta t}=20\left(n+\frac{1}{4}\right)\Ums$；若波沿 $x$ 轴负向传播，由题图知波传播的最短距离为 $\frac{3\lambda}{4}$，即传播的最短时间为 $\frac{3T}{4}$，由波的周期性知 $\Delta t=\left(n+\frac{3}{4}\right)T$，$n$ 取自然数，由 $v=\frac{\lambda}{T}$ 得 $v=\frac{4\left(n+\frac{3}{4}\right)}{\Delta t}=20\left(n+\frac{3}{4}\right)\Ums$。
}

\item
%% number: 24
%% paperName: 1989年全国高考第 24 题 3 分
%% typeId: 3
%% type: 填空题
%% chapter: 电路
%% point: 非纯电阻电路
%% method: 无
%% score: 3
%% degree: 550
%% duplicateId: 0
%% body:
某一用直流电动机提升重物的装置，如上右图所示。重物的质量 $m=50\Ukg$，电源的电动势 $E=110\UV$，不计电源内阻及各处的摩擦。当电动机以 $v=0.90\Ums$ 的恒定速度向上提升重物时，电路中的电流强度 $I=5\UA$，由此可知电动机线圈的电阻 $R=$\tkanswer[1.8cm]{$4.4\UO$}。
\onepicture[align=right, w=5cm]{\includesvg{figs/24}}

%% answer: 4.4（答 4.36 或 4.0 的亦给 3 分）
\jdanswer{
$4.4\UO$（答 $4.36\UO$ 或 $4.0\UO$ 的亦给 3 分）
}

%% memo:
\memoanswer{
由功率的表达式得 $P=Fv=mgv=50\times9.8\times0.9\UW=441\UW$，由电功率的表达式知，电机的输出功率为 $P=EI-I^{2}R$，联立解得 $R=4.36\UO\approx4.4\UO$。
}

\item
%% number: 25
%% paperName: 1989年全国高考第 25 题 3 分
%% typeId: 3
%% type: 填空题
%% chapter: 直线运动
%% point: 时间中点速度
%% method: 无
%% score: 3
%% degree: 650
%% duplicateId: 0
%% body:
在测定匀变速直线运动的加速度的实验中，用打点计时器记录纸带运动的时间。计时器所用电源的频率为 $50\UHz$。图为做匀变速直线运动的小车带动的纸带上记录的一些点，在每相邻的两点中间都有四个点未画出。按时间顺序取 0、1、2、3、4、5 六个点，用米尺量出 1、2、3、4、5 点到 0 点的距离分别是 8.78、16.08、21.87、26.16、28.94（单位：$\Ucm$），由此可得小车的加速度的大小为\tkanswer[1.8cm]{$1.50\Umsq$}，方向\tkanswer[2.2cm]{与速度方向相反}。
\onepicture[align=right, w=8cm]{figs/25.png}

%% answer: 1.50，与速度方向相反。
\jdanswer{
$1.50\Umsq$，与速度方向相反。
}

%% memo:
\memoanswer{
由题意知，相邻两点间的距离为 $s_{1}=8.78\Ucm$，$s_{2}=7.30\Ucm$，$s_{3}=5.79\Ucm$，$s_{4}=4.29\Ucm$，$s_{5}=2.78\Ucm$，然后算出几个加速度：$a_{1}=\frac{s_{3}-s_{1}}{2T^{2}}=-1.495\Umsq$，$a_{2}=\frac{s_{4}-s_{2}}{2T^{2}}=-1.505\Umsq$，$a_{3}=\frac{s_{5}-s_{3}}{2T^{2}}=-1.505\Umsq$，取平均值 $\bar{a}=\frac{1}{3}(a_{1}+a_{2}+a_{3})=1.502\Umsq$，方向向左。
}

\item
%% number: 26
%% paperName: 1989年全国高考第 26 题 3 分
%% typeId: 3
%% type: 填空题
%% chapter: 电磁感应
%% point: 线圈过有界磁场
%% method: 无
%% score: 3
%% degree: 650
%% duplicateId: 0
%% body:
电阻为 $R$ 的矩形导线框 abcd，边长 ab $=L$，ad $=h$，质量为 $m$，自某一高度自由落下，通过一匀强磁场，磁场方向垂直纸面向里，磁场区域的宽度为 $h$（如图所示）。若线框恰好以恒定速度通过磁场，线框中产生的焦耳热是\tkanswer[1.8cm]{$2mgh$}。（不考虑空气阻力）
\onepicture[align=right, w=4cm]{\includesvg{figs/26}}

%% answer: $2mgh$
\jdanswer{
$2mgh$
}

%% memo:
\memoanswer{
解法一：由功能关系得 $Q=\Delta E_{\text{p}}=2mgh$。

解法二：线框切割产生的电动势为 $E=BLv$，闭合回路中的电流为 $I=\frac{E}{R}$，线框受到的安培力为 $F_{\text{安}}=\frac{B^{2}L^{2}v}{R}$，由线框受力平衡得 $mg=F_{\text{安}}$，由焦耳定律得 $Q=I^{2}Rt$，联立解得 $Q=2mgh$。
}

\item
%% number: 27
%% paperName: 1989年全国高考第 27 题 3 分
%% typeId: 3
%% type: 填空题
%% chapter: 力物体的平衡
%% point: 平衡综合
%% method: 整体与隔离
%% score: 3
%% degree: 650
%% duplicateId: 0
%% body:
质量为 $m$ 的运动员站在质量为 $m/2$ 的均匀长板 AB 的中点，板位于水平地面上，可绕通过 B 点的水平轴转动，板的 A 端系有轻绳，轻绳的另一端绕过两个定滑轮后，握在运动员手中。当运动员用力拉绳时，滑轮两侧的绳都保持在竖直方向，如图所示。要使板的 A 端离开地面，运动员作用于绳的最小拉力是\tkanswer[2cm]{$\frac{1}{2}mg$}。
\onepicture[align=right, w=5cm]{figs/27.png}

%% answer: $\frac{1}{2}mg$
\jdanswer{
$\frac{1}{2}mg$
}

%% memo:
\memoanswer{
设板长为 $l$，板的 A 端恰好离开地面时，由力矩平衡得 $F\cdot l=N'\frac{l}{2}+\frac{mg}{2}\cdot\frac{l}{2}$，$N'$ 为人对板的压力，隔离 A：$F+N=mg$，$N$ 为板对人的支持力，联立解得 $F=\frac{1}{2}mg$。
}

\item
%% number: 28
%% paperName: 1989年全国高考第 28 题 3 分
%% typeId: 3
%% type: 填空题
%% chapter: 气体性质
%% point: 阿伏伽德罗常数
%% method: 无
%% score: 3
%% degree: 650
%% duplicateId: 0
%% body:
一个房间的地面面积是 $15\Umq$，高 $3\Um$。试估算该房间内空气的质量。已知空气的平均摩尔质量是 $2.9\times10^{-2}\Ukgmol$。答：\tkanswer[1.6cm]{$58\Ukg$}

%% answer: 58（答案在 50 到 60 之间的都给 3 分）
\jdanswer{
58（答案在 50 到 60 之间的都给 3 分）
}

%% memo:
\memoanswer{
标准情况下，$1\Umol$ 的空气体积为 $22.4\times10^{-3}\Umc$，所以房间内空气的摩尔数为 $n=\frac{15\times3}{22.4\times10^{-3}}=2.0\times10^{3}$，故 $m=nm_{0}=2.0\times10^{3}\times2.9\times10^{-2}\Ukg=58\Ukg$。
}

\item
%% number: 29
%% paperName: 1989年全国高考第 29 题 10 分
%% typeId: 4
%% type: 实验题
%% chapter: 电路
%% point: 电阻定律
%% method: 无
%% score: 10
%% degree: 550
%% duplicateId: 0
%% body:
在测定金属的电阻率的实验中，金属导线长约 0.8 米，直径小于 1 毫米，电阻在 5 欧左右。实验步骤如下：
\begin{enumerate}
	\item
	用米尺测量金属导线的长度，测三次，求出平均值 $L$。在金属导线三个不同的位置上用\tkanswer[2.5cm]{螺旋测微器（或千分尺）}测量直径，求出平均值 $d$。
	\item
	用伏安法测量金属导线的电阻 $R$。试把下图中所给的器材连接成测量 $R$ 的合适的线路。图中安培表的量程为 0.6 安，内阻接近 1 欧；伏特表的量程为 3 伏，内阻为几千欧；电源的电动势为 6 伏；变阻器的阻值为 $0\sim20$ 欧。在闭合电键前，变阻器的滑动触点应处于正确位置。
	\item
	用上面测得的金属导线长度 $L$、直径 $d$ 和电阻 $R$，可根据电阻率的表达式 $\rho=$\tkanswer[2.8cm]{$\frac{\pi d^{2}}{4L}R$} 算出所测金属的电阻率。
\end{enumerate}
\onepicture[align=right, w=7cm]{figs/29.png}

%% answer: （1）螺旋测微器（或千分尺）；（2）连线如图；（3）$\frac{\pi d^{2}}{4L}R$
\jdanswer{
\begin{enumerate}
	\item
	螺旋测微器（或千分尺）
	\item
	连线如图所示

\onepicture[w=7cm]{figs/29_an.png}
	\item
	$\frac{\pi d^{2}}{4L}R$
\end{enumerate}
}

%% memo:
\memoanswer{
由电表阻值与待测电阻的比值 $\frac{R_{x}}{R_{A}}\ll\frac{R_{V}}{R_{x}}$ 知，应采用电流表外接法，变阻器接成限流器，滑动变阻器置于待测电压最小位置。由电阻定律得 $R=\rho\frac{L}{S}$，又 $S=\frac{1}{4}\pi d^{2}$，解得 $\rho=\frac{\pi d^{2}R}{4L}$。
}

\item
%% number: 30
%% paperName: 1989年全国高考第 30 题 9 分
%% typeId: 6
%% type: 计算题
%% chapter: 电磁感应
%% point: 电磁感应受力分析
%% method: 无
%% score: 9
%% degree: 750
%% duplicateId: 0
%% body:
如图所示，AB、CD 是两根足够长的固定平行金属导轨，两导轨间的距离为 $L$，导轨平面与水平面的夹角是 $\theta$。在整个导轨平面内都有垂直于导轨平面斜向上方的匀强磁场，磁感应强度为 $B$。在导轨的 AC 端连接一个阻值为 $R$ 的电阻。一根垂直于导轨放置的金属棒 ab，质量为 $m$，从静止开始沿导轨下滑，求 ab 棒的最大速度。

要求画出 ab 棒的受力图。已知 ab 与导轨间的滑动摩擦系数 $\mu$，导轨和金属棒的电阻都不计。
\onepicture[align=right, w=5cm]{figs/30.png}

%% answer: $v=\frac{mg(\sin\theta-\mu\cos\theta)R}{B^{2}L^{2}}$
\jdanswer{
$v=\frac{mg(\sin\theta-\mu\cos\theta)R}{B^{2}L^{2}}$
}

%% memo:
\memoanswer{
在下滑过程中，金属棒 ab 受到重力 $mg$，支持力 $N=mg\cos\theta$，摩擦力 $f=\mu mg\cos\theta$ 和安培力作用，方向如图所示。

\onepicture[w=3.5cm]{figs/30_an.png}

金属棒切割产生的电动势为 $E=BLv$，回路中的电流大小为 $I=\frac{E}{R}$，金属棒受到的安培力大小为 $F=BIL=\frac{B^{2}L^{2}v}{R}$。金属棒速度达到最大时，金属棒达到平衡，由平衡条件得 $mg\sin\theta-\mu mg\cos\theta-\frac{B^{2}L^{2}v}{R}=0$，解得最大速度 $v=\frac{mg(\sin\theta-\mu\cos\theta)R}{B^{2}L^{2}}$。
}

\item
%% number: 31
%% paperName: 1989年全国高考第 31 题 8 分
%% typeId: 6
%% type: 计算题
%% chapter: 光学
%% point: 透镜
%% method: 无
%% score: 8
%% degree: 650
%% duplicateId: 0
%% body:
把一个点光源放在焦距为 $f$ 的凸透镜的焦点上，在透镜的另一侧 2 倍焦距处放一个垂直于主轴的光屏，在光屏上看到一个半径为 $R$ 的光亮的圆。现保持透镜和光屏不动，而在主轴上移动点光源，若要使光屏上亮圆的半径缩为 $R/2$，则这个点光源应移到什么位置上？

%% answer: $u=\frac{4}{3}f$ 或 $u'=4f$
\jdanswer{
$u=\frac{4}{3}f$ 或 $u'=4f$
}

%% memo:
\memoanswer{
出射光束平行于主轴，因此亮圆的半径 $R$ 就等于透镜的半径。与光屏上半径为 $R/2$ 的亮圆相应的像点有两个，像距 $v$ 和 $v'$ 分别满足关系 $\frac{v-2f}{\frac{R}{2}}=\frac{v}{R}$，$\frac{2f-v'}{\frac{R}{2}}=\frac{v'}{R}$，代入成像公式 $\frac{1}{u}+\frac{1}{v}=\frac{1}{f}$，就得到点光源与透镜光心的距离分别为 $u=\frac{4}{3}f$ 或 $u'=4f$。

\onepicture[w=8cm]{figs/31_an.png}
}

\item
%% number: 32
%% paperName: 1989年全国高考第 32 题 9 分
%% typeId: 6
%% type: 计算题
%% chapter: 电场
%% point: 带电粒子的往复运动
%% method: 无
%% score: 9
%% degree: 550
%% duplicateId: 0
%% body:
一个质量为 $m$、带有电荷 $-q$ 的小物体，可在水平轨道 $Ox$ 上运动，$O$ 端有一与轨道垂直的固定墙。轨道处于匀强电场中，场强大小为 $E$，方向沿 $Ox$ 轴正向，如图所示。小物体以初速 $v_{0}$ 从 $x_{0}$ 点沿 $Ox$ 轨道运动，运动时受到大小不变的摩擦力 $f$ 作用，且 $f<qE$；设小物体与墙碰撞时不损失机械能，且电量保持不变，求它在停止运动前所通过的总路程 $s$。
\onepicture[align=right, w=7cm]{\includesvg{figs/32}}

%% answer: $s=\frac{2qEx_{0}+mv_{0}^{2}}{2f}$
\jdanswer{
$s=\frac{2qEx_{0}+mv_{0}^{2}}{2f}$
}

%% memo:
\memoanswer{
小物体受到的电场力 $F=-qE$，大小不变，方向指向墙；摩擦力 $f$ 的方向与小物体运动方向相反。不管开始时小物体是沿 $x$ 轴正方向或负方向运动，小物体在多次与墙碰撞后，最后将停止在原点 O 处。设小物体在停止前所通过的总路程为 $s$，在这个过程中，电势能减少了 $\Delta E_{\text{p}}=qEx_{0}$，小物体动能减少了 $\frac{1}{2}mv_{0}^{2}$。由于小物体与墙碰撞时不损失机械能，因而小物体克服摩擦力所做的功就等于所减少的动能和电势能之和，有 $fs=\frac{1}{2}mv_{0}^{2}+qEx_{0}$，解得小物体在停止前所通过的总路程 $s$ 等于 $s=\frac{2qEx_{0}+mv_{0}^{2}}{2f}$。
}

\end{enumerate}

\end{document}
